% Generated draft: requires mathematical and source review.
\documentclass[11pt]{article}
\usepackage[margin=25mm]{geometry}
\usepackage{fontspec}
\setmainfont{Latin Modern Roman}
\usepackage{amsmath,amssymb,mathtools,unicode-math}
\setmathfont{Latin Modern Math}
\usepackage{enumitem,needspace,xurl,hyperref}
\setlength{\parindent}{0pt}
\setlength{\parskip}{5pt}
\setlength{\emergencystretch}{3em}
\newcommand{\N}{\mathbb{N}}
\newcommand{\Z}{\mathbb{Z}}
\newcommand{\Q}{\mathbb{Q}}
\newcommand{\R}{\mathbb{R}}
\newcommand{\C}{\mathbb{C}}
\newcommand{\F}{\mathbb{F}}
\newcommand{\A}{\mathbb{A}}
\newcommand{\PP}{\mathbb{P}}
\newcommand{\E}{\mathbb{E}}
\newcommand{\eps}{\varepsilon}
\newcommand{\dd}{\,\mathrm{d}}
\begin{document}
\section*{UnsolvedMath ID: 11000004}
\textbf{Source problem code:} catalogue entry
\subsection{Definitions and mathematical statement}
The extraction has lost the decorations distinguishing a space from its completion. Write \(Y_g\) for the original space and \(\overline Y_g\) for the proposed completion, solely to distinguish the two roles. The preceding definitions of these spaces, their construction, the allowed genera and the function \(L\) are not supplied.

The problem asks to complete that construction and prove that
\[
 L:Y_g\longrightarrow[0,\infty)
 \quad\text{extends to}\quad
 \overline L:\overline Y_g\longrightarrow[0,\infty)
\]
as a proper map, meaning continuous with compact inverse image of every compact set. Ideally the inclusion \(Y_g\hookrightarrow\overline Y_g\) should be a homotopy equivalence, and \(\overline Y_g\setminus Y_g\) should be the boundary of a manifold-with-corners structure on \(\overline Y_g\). These latter goals are explicitly aspirational in the source.

Study \(\overline L\) as a ``Morse'' function, allowing nonisolated critical points, and explain the change in homotopy type of
\(\overline L^{-1}([0,t])\) when \(t\) passes through a critical value. The required smooth/stratified structure and meaning of critical point are not defined in the excerpt.

For \(g=2\), the source describes added pieces with extracted label \(Y_1^1\times Y_1^1\times S^1\), identified there with \(\mathbb H^2\times\mathbb H^2\times S^1\). The restricted function factors through the first two factors, with a unique minimum and no other critical points, hence a circle of minima. Its stated hope is that a point of \(\overline Y_2\) is a local minimum exactly when it is in one of these minimum circles, and that additional isolated index-\(1\) critical points connect those circles in a tree-like pattern. This would give \(\overline Y_2\) and \(Y_2\) the homotopy type of a wedge with one circle for every intended splitting, recovering the cited Mess theorem. The label convention and the indexing splittings are missing, and the proposed critical-point picture is not asserted as proved.
\subsection{Short English statement}
Construct the proposed completion of \(Y_g\), extend its length-like function properly, and analyze sublevel-set topology, with a conjectural genus-two minimum-circle and index-one critical-point picture recovering Mess's result.
\subsection{Sources}
\begin{itemize}
\item Farb, chapter in \emph{Problems on Mapping Class Groups and Related Topics}, Problem 2.3, supplied PDF page 13. \url{https://www.math.uchicago.edu/~farb/papers/mcgbook.pdf}. Recorded access 2026-07-29; extraction pdf-font-structure. The preceding completion construction and definitions are absent.
\item The supplied summary cites Mess (1992) for the theorem that \(I_2\) is free of infinite rank, without a full bibliography entry. The source says the hoped-for genus-two analysis should recover the previously mentioned theorem.
\item Local UnsolvedMath registry, \texttt{data/unsolvedmath/problems.json}, permanent ID \(11000004\), exact source code \texttt{catalogue entry}; supplied OPEN-TRIAGE status and programmatic open report.
\end{itemize}
\end{document}
